% QUESTAO
Após determinar o raio da círculo abaixo, a sua área, em metros quadrados, fica:

\parbox{0.4\linewidth}{

% ALTERNATIVAS
$52\pi$ m$^2$
$51\pi$ m$^2$
$53\pi$ m$^2$
$54\pi$ m$^2$
$55\pi$ m$^2$

}\parbox{0.59\linewidth}{
\begin{center}
\vspace{-6mm}
	\begin{tikzpicture}[>=latex,scale=0.2,rotate=0]
	\pgfmathsetmacro{\R}{sqrt(52)};
	\coordinate (O) at (0,0);   \fill (O) circle (0.4*\R mm);
	\draw[thick] (O) circle (\R);
	\draw[thick,nodelabel={-10pt}{$R$}] (O) -- ($(-45:4)+(-135:6)$);
	
	\draw[thick,nodelabel={10pt}{$4$\,m}] (O) -- (-45:4);
	
	\draw[thick] ($(-45:4)+(-135:6)$) -- ($(-45:4)+(45:6)$);
	\node at ($($(-45:5.5)+(-135:6)$)!0.6!(-45:5.5)$) {$6$\,m};
	
	\draw[dashed] ($(-45:4)+(-135:6)$) -- ($(-45:4)+(-135:6)+(-45:4.2)$) ($(-45:4)+(45:6)$) -- ($(-45:4)+(45:6)+(-45:4.2)$);
	
	\draw[ decorate, decoration={brace,raise=0pt} ]  ($(-45:4)+(45:6)+(-45:4.2)$) -- ($(-45:4)+(-135:6)+(-45:4.2)$);
	
	\node at (-45:10.5) {$12$\,m};
	
	\draw[rotate=45] (-90:4) rectangle +(-1,1);
	\fill[rotate=45] ($(-90:4)+(-0.5,0.5)$) circle (0.1);
	
	\end{tikzpicture}
\vspace{-9mm}
\end{center}
}


